\documentclass[11pt,letterpaper]{article}

\usepackage[utf8]{inputenc}
\usepackage[T1]{fontenc}
\usepackage[spanish]{babel}
\usepackage[margin=2.5cm]{geometry}
\usepackage{graphicx}
\usepackage{hyperref}
\usepackage{listings}
\usepackage{xcolor}
\usepackage{titlesec}
\usepackage{enumitem}
\usepackage{tcolorbox}

% Colores para código
\definecolor{codebg}{RGB}{245,240,230}
\definecolor{codecomment}{RGB}{100,120,100}
\definecolor{codestring}{RGB}{180,60,60}
\definecolor{codekw}{RGB}{60,60,150}

\lstset{
  backgroundcolor=\color{codebg},
  basicstyle=\ttfamily\small,
  commentstyle=\color{codecomment}\itshape,
  keywordstyle=\color{codekw}\bfseries,
  stringstyle=\color{codestring},
  breaklines=true,
  frame=single,
  rulecolor=\color{gray},
  showstringspaces=false,
}

\title{\Huge \textbf{TalkVinyl.} \\ \Large Marketplace de Vinilos con Arquitectura Hexagonal}
\author{Illyana Julienne Engle Reyes}
\date{\today}

\begin{document}

\maketitle

\tableofcontents
\newpage

\section{Introducción}

\textbf{TalkVynil} es una plataforma de comercio electrónico para venta de vinilos de música. El sistema implementa una \textbf{arquitectura hexagonal} (Ports \& Adapters) con separación estricta de capas, permitiendo que la lógica de negocio sea independiente de los frameworks y las bases de datos.

\subsection{Objetivos}

\begin{itemize}
  \item Modelar tres entidades interconectadas: \textbf{Usuario}, \textbf{Producto} y \textbf{Pedido}.
  \item Implementar CRUDs completos para cada entidad con control de acceso por roles.
  \item Asegurar las credenciales mediante hashing con bcrypt.
  \item Gestionar imágenes de productos fuera de la base de datos.
  \item Aplicar arquitectura hexagonal en el backend.
  \item Construir una SPA en React con módulos por rol.
\end{itemize}

\section{Arquitectura del Sistema}

\subsection{Capas}

El backend está organizado en cuatro capas:

\begin{enumerate}
  \item \textbf{Dominio:} entidades puras (\texttt{User}, \texttt{Product}, \texttt{Order}) y puertos de persistencia.
  \item \textbf{Aplicación:} servicios con la lógica de los casos de uso.
  \item \textbf{Infraestructura:} adaptadores que implementan los puertos (PostgreSQL, bcrypt, multer, svg-captcha).
  \item \textbf{Interfaces:} controladores HTTP (Express) y middlewares de autenticación y roles.
\end{enumerate}

\subsection{Diagrama}

\begin{verbatim}
┌─────────────────────────────────────┐
│         INTERFACES (HTTP)           │
│    Controllers + Middlewares        │
└──────────────┬──────────────────────┘
               │
┌──────────────▼──────────────────────┐
│      APPLICATION (Servicios)        │
│   userService / productService      │
│   orderService / captchaStore       │
└──────────────┬──────────────────────┘
               │
┌──────────────▼──────────────────────┐
│            DOMAIN                   │
│   Entidades + Puertos + Reglas      │
└──────────────▲──────────────────────┘
               │
┌──────────────┴──────────────────────┐
│        INFRASTRUCTURE               │
│   PostgreSQL / bcrypt / multer      │
└─────────────────────────────────────┘
\end{verbatim}

\section{Modelo de Base de Datos}

\subsection{Tabla \texttt{usuarios}}

\begin{itemize}
  \item \texttt{id} SERIAL PRIMARY KEY
  \item \texttt{nombre} VARCHAR(100)
  \item \texttt{email} VARCHAR(150) UNIQUE
  \item \texttt{password\_hash} VARCHAR(255)
  \item \texttt{rol} VARCHAR(20) — \texttt{admin}, \texttt{producto}, \texttt{pedido} o NULL
  \item \texttt{estado} VARCHAR(20) — \texttt{pendiente}, \texttt{activo}, \texttt{rechazado}
  \item \texttt{created\_at} TIMESTAMP
\end{itemize}

\subsection{Tabla \texttt{productos}}

\begin{itemize}
  \item \texttt{id} SERIAL PRIMARY KEY
  \item \texttt{nombre}, \texttt{artista}, \texttt{descripcion}
  \item \texttt{precio} NUMERIC(10,2)
  \item \texttt{imagen\_url} VARCHAR(255) — ruta al archivo (no el binario)
  \item \texttt{stock} INT
  \item \texttt{estado} VARCHAR(20) — \texttt{pendiente}, \texttt{aprobado}, \texttt{rechazado}, \texttt{vendido}
  \item \texttt{creado\_por} INT REFERENCES usuarios(id)
\end{itemize}

\subsection{Tabla \texttt{pedidos}}

\begin{itemize}
  \item \texttt{id} SERIAL PRIMARY KEY
  \item \texttt{usuario\_id}, \texttt{producto\_id}
  \item \texttt{cantidad} INT
  \item \texttt{total} NUMERIC(10,2)
  \item \texttt{estado} VARCHAR(20)
\end{itemize}

\section{Endpoints Principales}

\subsection{Autenticación}
\begin{itemize}
  \item \texttt{POST /usuarios} — Registro con captcha
  \item \texttt{POST /login} — Login con captcha → JWT
  \item \texttt{GET /captcha} — Genera captcha SVG
\end{itemize}

\subsection{Usuarios (admin)}
\begin{itemize}
  \item \texttt{GET /usuarios} — Listar
  \item \texttt{PUT /usuarios/:id/aprobar} — Aprobar y asignar rol
  \item \texttt{PUT /usuarios/:id/rol} — Cambiar rol
  \item \texttt{DELETE /usuarios/:id} — Eliminar
\end{itemize}

\subsection{Productos}
\begin{itemize}
  \item \texttt{GET /productos} — Listar según rol
  \item \texttt{GET /productos/catalogo} — Catálogo público
  \item \texttt{POST /productos} — Crear (multipart con imagen)
  \item \texttt{PUT /productos/:id/aprobar} — Aprobar
  \item \texttt{DELETE /productos/:id} — Eliminar
\end{itemize}

\subsection{Pedidos}
\begin{itemize}
  \item \texttt{POST /pedidos} — Crear pedido
  \item \texttt{GET /pedidos} — Listar según rol
  \item \texttt{PUT /pedidos/:id/aprobar} — Aprobar pedido
  \item \texttt{GET /pedidos/mis-solicitudes} — Solicitudes del vendedor
\end{itemize}

\section{Guía de Despliegue en WSL}

\subsection{Requisitos}
\begin{itemize}
  \item WSL (Ubuntu)
  \item Node.js 20+
  \item PostgreSQL 14+
  \item Git
\end{itemize}

\subsection{Pasos}

\begin{lstlisting}[language=bash]
# 1. Clonar
git clone https://github.com/illyanaengle22-ui/Node-JS-Mysql.git webapp
cd webapp

# 2. Base de datos
sudo -u postgres psql -c "CREATE DATABASE practica_api;"
sudo -u postgres psql -d practica_api -f backend/migrations/001_init_vinilos.sql

# 3. Backend
cd backend && npm install
# Crear src/.env con DB_HOST, DB_USER, DB_PASSWORD, DB_NAME, JWT_SECRET
node src/server.js

# 4. Frontend (nueva terminal)
cd ../frontend && npm install
npm run dev
\end{lstlisting}

\subsection{Credenciales de prueba}

\begin{center}
\begin{tabular}{|l|l|}
\hline
\textbf{Email} & \textbf{Contraseña} \\
\hline
admin@vinilos.com & Admin123! \\
\hline
\end{tabular}
\end{center}

\section{Conclusiones}

El proyecto demuestra la aplicación práctica de arquitectura hexagonal en un caso real de comercio electrónico, con separación estricta de responsabilidades, seguridad de credenciales y control de acceso por roles. La solución es mantenible, escalable y testeable.

\end{document}
